\documentclass{amsart}
% For dealing with references we use the comment environment
\usepackage{verbatim}
\newenvironment{reference}{\comment}{\endcomment}
%\newenvironment{reference}{}{}
\newenvironment{slogan}{\comment}{\endcomment}
\newenvironment{history}{\comment}{\endcomment}

% For commutative diagrams we use Xy-pic
\usepackage[all]{xy}

% We use 2cell for 2-commutative diagrams.
\xyoption{2cell}
\UseAllTwocells

% We use multicol for the list of chapters between chapters
\usepackage{multicol}

% This is generally recommended for better output
\usepackage{lmodern}
\usepackage[T1]{fontenc}

% For cross-file-references

% Package for hypertext links:
\usepackage{hyperref}

% For any local file, say "hello.tex" you want to link to please

% Theorem environments.
%
\theoremstyle{plain}
\newtheorem{theorem}[subsection]{Theorem}
\newtheorem{proposition}[subsection]{Proposition}
\newtheorem{lemma}[subsection]{Lemma}

\theoremstyle{definition}
\newtheorem{definition}[subsection]{Definition}
\newtheorem{example}[subsection]{Example}
\newtheorem{exercise}[subsection]{Exercise}
\newtheorem{situation}[subsection]{Situation}

\theoremstyle{remark}
\newtheorem{remark}[subsection]{Remark}
\newtheorem{remarks}[subsection]{Remarks}

\numberwithin{equation}{subsection}

% Macros
%
\def\lim{\mathop{\mathrm{lim}}\nolimits}
\def\colim{\mathop{\mathrm{colim}}\nolimits}
\def\Spec{\mathop{\mathrm{Spec}}}
\def\Hom{\mathop{\mathrm{Hom}}\nolimits}
\def\Ext{\mathop{\mathrm{Ext}}\nolimits}
\def\SheafHom{\mathop{\mathcal{H}\!\mathit{om}}\nolimits}
\def\SheafExt{\mathop{\mathcal{E}\!\mathit{xt}}\nolimits}
\def\Sch{\mathit{Sch}}
\def\Mor{\mathop{\mathrm{Mor}}\nolimits}
\def\Ob{\mathop{\mathrm{Ob}}\nolimits}
\def\Sh{\mathop{\mathit{Sh}}\nolimits}
\def\NL{\mathop{N\!L}\nolimits}
\def\CH{\mathop{\mathrm{CH}}\nolimits}
\def\proetale{{pro\text{-}\acute{e}tale}}
\def\etale{{\acute{e}tale}}
\def\QCoh{\mathit{QCoh}}
\def\Ker{\mathop{\mathrm{Ker}}}
\def\Im{\mathop{\mathrm{Im}}}
\def\Coker{\mathop{\mathrm{Coker}}}
\def\Coim{\mathop{\mathrm{Coim}}}

% Boxtimes
%
\DeclareMathSymbol{\boxtimes}{\mathbin}{AMSa}{"02}

%
% Macros for moduli stacks/spaces
%
\def\QCohstack{\mathcal{QC}\!\mathit{oh}}
\def\Cohstack{\mathcal{C}\!\mathit{oh}}
\def\Spacesstack{\mathcal{S}\!\mathit{paces}}
\def\Quotfunctor{\mathrm{Quot}}
\def\Hilbfunctor{\mathrm{Hilb}}
\def\Curvesstack{\mathcal{C}\!\mathit{urves}}
\def\Polarizedstack{\mathcal{P}\!\mathit{olarized}}
\def\Complexesstack{\mathcal{C}\!\mathit{omplexes}}
% \Pic is the operator that assigns to X its picard group, usage \Pic(X)
% \Picardstack_{X/B} denotes the Picard stack of X over B
% \Picardfunctor_{X/B} denotes the Picard functor of X over B
\def\Pic{\mathop{\mathrm{Pic}}\nolimits}
\def\Picardstack{\mathcal{P}\!\mathit{ic}}
\def\Picardfunctor{\mathrm{Pic}}
\def\Deformationcategory{\mathcal{D}\!\mathit{ef}}

\setlength{\emergencystretch}{3em}
\begin{document}
\title[Stacks Project, Tag 0CBX]{573 --- A flat finite-type family over a Noetherian local base has completed equation xy=h at a split node of its fibre}
\maketitle
\noindent Copyright (C) 2005--2025 Johan de Jong. Stacks Project authors. Source: \href{https://stacks.math.columbia.edu/tag/0CBX}{Tag 0CBX}.

\noindent Editorial difficulty note (not author text). Difficulty H2: The proof constructs compatible lifts of the two branch parameters by successive corrections in every power of the base maximal ideal. Completeness produces an exact equation, and flatness plus the special-fibre isomorphism is used to prove the completed algebra map is both surjective and injective..

\noindent Editorial provenance note (not author text). History: excerpt from revision \texttt{a0444\allowbreak 6e57e\allowbreak c1fbc\allowbreak 252a8\allowbreak 71afc\allowbreak ec775\allowbreak 2fb28\allowbreak 07b14}, curves.tex, lines 4919--5015. Reference presentation was adapted; original-author context and core support blocks were retained or added. The mathematical wording is unchanged. Licensed under GNU FDL 1.2 or later; no invariant sections or cover texts. See ../licenses/COPYING. Context blocks reproduce author definitions and supporting results; they are not additional counted problems.

\begin{definition}
\label{definition-split-node}
Let $k$ be a field. Let $X$ be a $1$-dimensional algebraic $k$-scheme.
Let $x \in X$ be a closed point. We say $x$ is a {\it split node}
if $x$ is a node, $\kappa(x) = k$, and the equivalent assertions of
Remark \ref{remark-trivial-quadratic-extension}
hold for $A = \mathcal{O}_{X, x}$.
\end{definition}

\begin{remark}[Trivial quadratic extension]
\label{remark-trivial-quadratic-extension}
Let $k$ be a field. Let $(A, \mathfrak m, \kappa)$ be as in
Remark \href{https://stacks.math.columbia.edu/tag/0CBT}{remark quadratic extension (Tag 0CBT)} and let $\kappa'/\kappa$
be the associated separable algebra of degree $2$.
Then the following are equivalent
\begin{enumerate}
\item $\kappa' \cong \kappa \times \kappa$ as $\kappa$-algebra,
\item the form $q$ of Lemma \href{https://stacks.math.columbia.edu/tag/0C49}{lemma nodal algebraic (Tag 0C49)}
can be chosen to be $xy$,
\item $A$ has two branches,
\item the extension $A'/A$ of Lemma \href{https://stacks.math.columbia.edu/tag/0C4A}{lemma 2 branches delta 1 (Tag 0C4A)}
has two maximal ideals, and
\item $A^\wedge \cong \kappa[[x, y]]/(xy)$ as a $k$-algebra.
\end{enumerate}
The equivalence between these conditions has been shown in the
proof of Lemma \href{https://stacks.math.columbia.edu/tag/0C4A}{lemma 2 branches delta 1 (Tag 0C4A)}. If $X$ is a
locally Noetherian $k$-scheme and $x \in X$ is a point such that
$\mathcal{O}_{X, x} = A$, then this means exactly that
there are two points $x_1, x_2$ of the normalization $X^\nu$
lying over $x$ and that $\kappa(x) = \kappa(x_1) = \kappa(x_2)$.
\end{remark}

\begin{lemma}
\label{lemma-formal-local-structure-nodal-family}
Let $f : T \to S$ be a morphism of schemes. Let $t \in T$
with image $s \in S$. Assume
\begin{enumerate}
\item $f$ is flat at $t$,
\item $\mathcal{O}_{S, s}$ is Noetherian,
\item $f$ is locally of finite type,
\item $t$ is a split node of the fibre $T_s$.
\end{enumerate}
Then there exists an $h \in \mathfrak m_s^\wedge$ and an isomorphism
$$
\mathcal{O}_{T, t}^\wedge \cong
\mathcal{O}_{S, s}^\wedge[[x, y]]/(xy - h)
$$
of $\mathcal{O}_{S, s}^\wedge$-algebras.
\end{lemma}

\begin{proof}
We replace $S$ by $\Spec(\mathcal{O}_{S, s})$ and $T$ by the base change
to $\Spec(\mathcal{O}_{S, s})$. Then $T$ is locally Noetherian and hence
$\mathcal{O}_{T, t}$ is Noetherian.
Set $A = \mathcal{O}_{S, s}^\wedge$, $\mathfrak m = \mathfrak m_A$, and
$B = \mathcal{O}_{T, t}^\wedge$. By
More on Algebra, Lemma \href{https://stacks.math.columbia.edu/tag/0C4G}{lemma flat completion (Tag 0C4G)}
we see that $A \to B$ is flat. Since
$\mathcal{O}_{T, t}/\mathfrak m_s \mathcal{O}_{T, t} = \mathcal{O}_{T_s, t}$
we see that $B/\mathfrak m B = \mathcal{O}_{T_s, t}^\wedge$.
By assumption (4) and Lemma \href{https://stacks.math.columbia.edu/tag/0CBW}{lemma split node (Tag 0CBW)}
we conclude there exist $\overline{u}, \overline{v} \in B/\mathfrak m B$
such that the map
$$
(A/\mathfrak m)[[x, y]] \longrightarrow B/\mathfrak m B,\quad
x \longmapsto \overline{u},
x \longmapsto \overline{v}
$$
is surjective with kernel $(xy)$.

\medskip\noindent
Assume we have $n \geq 1$ and $u, v \in B$ mapping to
$\overline{u}, \overline{v}$ such that
$$
u v = h + \delta
$$
for some $h \in A$ and $\delta \in \mathfrak m^nB$.
We claim that there exist $u', v' \in B$ with
$u - u', v - v' \in \mathfrak m^n B$ such that
$$
u' v' = h' + \delta'
$$
for some $h' \in A$ and $\delta' \in \mathfrak m^{n + 1}B$.
To see this, write $\delta = \sum f_i b_i$ with
$f_i \in \mathfrak m^n$ and $b_i \in B$. Then write
$b_i = a_i + u b_{i, 1} + v b_{i, 2} + \delta_i$ with
$a_i \in A$, $b_{i, 1}, b_{i, 2} \in B$ and $\delta_i \in \mathfrak m B$.
This is possible because the residue field of $B$ agrees with the
residue field of $A$ and the images of $u$ and $v$ in $B/\mathfrak m B$
generate the maximal ideal. Then we set
$$
u' = u - \sum b_{i, 2}f_i,\quad
v' = v - \sum b_{i, 1}f_i
$$
and we obtain
$$
u'v' = h + \delta - \sum (b_{i, 1}u + b_{i, 2}v)f_i + \sum
c_{ij}f_if_j =
h + \sum a_if_i + \sum f_i \delta_i + \sum c_{ij}f_if_j
$$
for some $c_{i, j} \in B$.
Thus we get a formula as above with $h' = h + \sum a_if_i$
and $\delta' = \sum f_i \delta_i + \sum c_{ij}f_if_j$.

\medskip\noindent
Arguing by induction and starting with any lifts $u_1, v_1 \in B$
of $\overline{u}, \overline{v}$ the result of the previous paragraph
shows that we find a sequence of elements
$u_n, v_n \in B$ and $h_n \in A$ such that
$u_n - u_{n + 1} \in \mathfrak m^n B$,
$v_n - v_{n + 1} \in \mathfrak m^n B$,
$h_n - h_{n + 1} \in \mathfrak m^n$,
and such that $u_n v_n - h_n \in \mathfrak m^n B$.
Since $A$ and $B$ are complete we can set
$u_\infty = \lim u_n$, $v_\infty = \lim v_n$,  and
$h_\infty = \lim h_n$, and then we obtain $u_\infty v_\infty = h_\infty$
in $B$. Thus we have an $A$-algebra map
$$
A[[x, y]]/(xy - h_\infty) \longrightarrow B
$$
sending $x$ to $u_\infty$ and $v$ to $v_\infty$.
This is a map of flat $A$-algebras which is an
isomorphism after dividing by $\mathfrak m$.
It is surjective modulo $\mathfrak m$ and hence surjective
by completeness and
Algebra, Lemma \href{https://stacks.math.columbia.edu/tag/0315}{lemma completion generalities (Tag 0315)}.
Then we can apply Algebra, Lemma \href{https://stacks.math.columbia.edu/tag/00ME}{lemma mod injective (Tag 00ME)}
to conclude it is an isomorphism.
\end{proof}
\end{document}
